\chapter{Conclusiones}
\label{cap:conclusiones}

El objetivo de este informe fue modelar BeeTracer con el paradigma estructurado, dando
continuidad a la reconstrucción por ingeniería inversa iniciada en el Informe 2. El análisis
partió del entorno del sistema: la declaración de propósitos fijó su razón de ser, registrar de
forma verificable el estado de la carga en cada etapa del desconsolidado, y el diagrama de
contexto y la lista de acontecimientos delimitaron los seis terminadores y los veinticuatro
estímulos a los que debe responder. Desde allí, los diagramas de flujo de datos descompusieron
el sistema en seis procesos principales, veintiocho procesos de nivel~1 y diez de nivel~2, cuyos
flujos y almacenes quedaron definidos sin ambigüedad en el diccionario de datos y cuya lógica
quedó descrita en las mini especificaciones.

Comparar ambos enfoques sobre el mismo sistema fue la principal lección del trabajo. El modelo
orientado a objetos del Informe 2 organizaba el sistema en torno a \emph{quién} es responsable
de cada dato y comportamiento; el modelo estructurado lo organiza en torno a \emph{qué} se hace
con los datos y dónde quedan. Este cambio de perspectiva hizo visibles detalles que el
diagrama de clases no mostraba con claridad: la necesidad de registrar la etapa de cada
fotografía para ordenar el reporte, la conveniencia de registrar cada envío del reporte como
un hecho propio y la obligación de registrar el contenedor cuando se opera en tarja ciega. A
la vez, la trazabilidad entre ambos modelos se mantuvo: las entidades del DER coinciden con las
clases de entidad, y los módulos predefinidos del diseño, con las clases de servicio.

La distinción entre el modelo esencial y el diseño resultó especialmente útil en BeeTracer.
Al suponer tecnología perfecta, el análisis se concentró en las reglas de negocio (evidencia
obligatoria, contraste de cantidades, validación del supervisor) sin distraerse con la falta
de señal en el patio. En el diseño, esos aspectos tecnológicos se incorporaron como módulos
bien delimitados (compresión, sellado, encolado y sincronización de fotografías), y el
análisis de transacción permitió obtener una arquitectura en la que cada subsistema
corresponde a un proceso del nivel~0. El módulo \este{Registrar evidencia fotográfica},
reutilizado por todas las etapas, muestra cómo un buen diseño concentra en un solo lugar la
regla que da valor al producto.

Como limitación, al igual que en el informe anterior, no se tuvo acceso al código ni a la
base de datos de BeeTracer, por lo que los procesos, módulos y tablas propuestos son una
reconstrucción coherente con lo observado y no necesariamente idéntica a la implementación
real. Aun así, el modelo relacional y las especificaciones de módulos constituyen una base
concreta para una eventual implementación o para evaluar mejoras del sistema, como la
auditoría de los cambios de estado o la integración con los módulos complementarios que la
empresa comercializa por separado.
